% =====================================================================
%  Section 6  --  Long-form lecture notes
%  Agentic AI: What Works, What Breaks, and What You Can't Trust
% =====================================================================
\documentclass[11pt]{article}

\usepackage[margin=1in]{geometry}
\usepackage[utf8]{inputenc}
\usepackage[T1]{fontenc}
\usepackage{lmodern}
\usepackage{microtype}
\usepackage{parskip}
\usepackage{hyperref}
\usepackage{natbib}
\usepackage{booktabs}
\usepackage{array}
\usepackage{enumitem}
\usepackage[dvipsnames]{xcolor}

\hypersetup{
  colorlinks=true,
  linkcolor=black,
  citecolor=black,
  urlcolor=blue!60!black
}

\title{Agentic AI: What Works, What Breaks, and What You Can't Trust\\
  {\large A peer-reviewed reality check for decision-makers}}
\author{Section 6 lecture notes}
\date{}

\begin{document}
\maketitle

\section*{Why another reality check, right now}

Section~5 made a specific claim, grounded in the largest AI-researcher
survey ever conducted \citep{grace2024thousands}: expert timelines for
human-level machine intelligence are unstable, narrow-task forecasts
already undershot, and the tail of catastrophic outcomes is not a
fringe position inside the research community. That lecture ended on
a question: what do you, a manager, actually do about it.

Section~6 picks up where that question lands. The honest answer is
not another strategic framework. The honest answer is that the
category managers are being sold today, ``agentic AI,'' behaves
differently from the chatbots their pilots trained them to evaluate.
It has scaffolding you can't see, benchmark scores you can't trust,
and failure modes that are well-documented in peer-reviewed literature
but have not yet made it into most procurement templates. This
section unpacks each of those three, and ends with a short playbook
you can take to Monday's meeting.

Everything below is anchored to papers published in venues where peer
review is real: ICML, NeurIPS, ICLR, FAccT, and \emph{Science} /
\emph{Patterns} (Cell Press). Where we cite a preprint or technical
report, we flag it explicitly. Agentic AI is a fast-moving field, and
a lot of the material that circulates in the trade press is not
peer-reviewed. Managers who treat every widely-shared finding as
settled science will be surprised, repeatedly.

\section{What an agent actually is}

The word ``agent'' has collapsed in meaning. In the research
literature and, increasingly, in procurement documents, it refers to
a loop that wraps a language model in three things: a \emph{planner}
that decides what to try, a set of \emph{tools} (a browser, a code
interpreter, a retrieval index, a company API), and a \emph{reflection}
mechanism that looks at the result and decides whether to retry.

A manager's first instinct is to ask which model is best. That is the
wrong lever. The single most replicated finding in the agentic-AI
literature is that the scaffolding, not the base model, accounts for
most of the performance gain.

\citet{yao2023react} introduced ReAct (ICLR 2023), which interleaves
a short reasoning step with a single tool action, then observes the
result before planning the next step. On ALFWorld, a household-task
benchmark, ReAct raised the same base model's success rate from
roughly 34\% to 71\%. The base model did not change; the loop around
it did.

\citet{shinn2023reflexion} extended the idea in Reflexion (NeurIPS
2023). After a failure, the agent writes a paragraph in plain
language about \emph{why} it failed, stores that paragraph, and
consults it on the next attempt. This simple trick raised GPT-4's
pass@1 on the HumanEval coding benchmark from roughly 80\% to 91\%.

\citet{schick2023toolformer} made the sharpest point (NeurIPS 2023):
a 6.7-billion-parameter model that teaches itself, in a self-supervised
way, to call external APIs outperforms GPT-3 (175 billion parameters)
on downstream tasks where tool use matters. Capability and model size
are not the same lever. If you are evaluating vendors on raw model
size, you are evaluating the wrong thing.

\citet{park2023generative}, at UIST 2023 (best-paper honorable
mention), gave 25 language-model agents memory, reflection, and
planning, then dropped them into a simulated town. Left alone, they
organised a Valentine's Day party. The point is not that a simulated
village is commercially important. The point is that the same
scaffolding ideas, applied at population scale, produce emergent
collective behaviour that no one explicitly programmed. Procurement
conversations should ask what happens when \emph{two} of your agents
interact, not just when one does.

The procurement implication is the durable lesson from this block of
research. You are not buying a model. You are buying a model plus
planner plus tools plus retry loop plus guardrails plus integration.
Vendor demos that highlight the model and hide the scaffolding are
demonstrating the wrong thing.

\section{What the benchmarks say, and what they hide}

The agentic-AI literature settled on four benchmarks, all peer-reviewed
at ICLR 2024: GAIA, SWE-bench, AgentBench, and WebArena. All four were
designed to measure something closer to real work than prior LLM
benchmarks like MMLU. All four are now being misquoted in vendor
pitches. This section lays out what they actually say.

\subsection*{GAIA (Mialon et al., ICLR 2024)}

GAIA is a set of 466 questions that a careful human can answer in
minutes but that require an agent to plan, browse, reason over files,
and use tools in combination. Questions are stratified into three
levels of difficulty. Humans score about 92\% overall. At release,
the best LLM-plus-plugins configuration reached roughly 15\%
\citep{mialon2024gaia}.

As of April 2026, the top score on the public GAIA leaderboard is
44.8\% (GPT-5 Mini). Humans remain at roughly 92\%. Progress in two
years: from 15\% to 44.8\%, a tripling. Gap remaining: 47 percentage
points. Rate of closure: fast. Honest reading: capability has grown
substantially; the gap is still the story.

\subsection*{SWE-bench (Jimenez et al., ICLR 2024, oral)}

SWE-bench is a set of 2,294 real GitHub issues from popular
open-source Python repositories, paired with the human-written patches
that actually fixed them \citep{jimenez2024swebench}. An agent is
given the issue text and the repository, and must produce a patch
that passes the project's own test suite. At release in October 2023,
GPT-4 resolved roughly 2\% of the issues.

As of April 2026, the top SWE-bench Verified score is 93.9\% (Claude
Mythos Preview). That is a forty-fold improvement in 30 months. But
there is a catch. In mid-2024, OpenAI stopped reporting Verified
scores for its own frontier models after detecting that every frontier
system they tested showed signs of training-data contamination on the
benchmark. This is documented on OpenAI's own announcement page for
the Verified subset \citep{openai2024swebenchverified}, and echoed in
independent audits such as the Berkeley RDI group's widely-read 2026
post ``How We Broke Top AI Agent Benchmarks.''

The managerial takeaway is uncomfortable but important: even when a
benchmark is peer-reviewed, human-audited, and widely cited, its
numbers can become a lagging indicator of capability and a leading
indicator of contamination, at the same time, and the two are
extremely difficult to separate from the outside.

\subsection*{WebArena (Zhou et al., ICLR 2024)}

WebArena gives an agent a browser and a realistic web task, such as
booking a flight or filing a reimbursement, inside a sandboxed copy
of real web software \citep{zhou2024webarena}. At release, GPT-4
succeeded on about 14\% of tasks end-to-end. As of January 2026, the
published state of the art is 71.6\% (OpAgent). The measured human
baseline on these tasks is approximately 78\%.

This is the most optimistic of the four benchmarks for deployment.
The gap to human baseline has shrunk to roughly seven percentage
points on at least one reported system, though contamination caveats
apply here too.

\subsection*{AgentBench (Liu et al., ICLR 2024)}

AgentBench \citep{liu2024agentbench} evaluates 25 language models
across 8 interactive environments ranging from OS command lines to
database queries to card games. Its headline finding from release is
less about any single number and more about a gap: the best
commercial models significantly outperform the best open-source
models across most environments. That gap has narrowed in the two
years since, but AgentBench scores are not tracked as aggressively as
SWE-bench or GAIA, so they are less susceptible to the contamination
dynamics that beset those two.

\subsection*{The benchmark crisis, in one paragraph}

Taken together, the four benchmarks tell a coherent story. Capability
is growing fast, in some cases faster than human improvement on the
same tasks. But the numbers you see in the trade press are unstable
in both directions. Genuine capability growth pushes them up.
Training-data contamination pushes them up as well. From outside the
lab that ran the test, these two are difficult to separate. A
procurement process that treats a quoted benchmark score as a trusted
signal is, in 2026, operating on an unreliable input.

\section{How agentic AI actually breaks}

Section~4 of this course catalogued the dark side of LLMs. Section~6
picks up the subset of those harms that manifest specifically when a
model is given autonomy, tools, or long horizons. Four failure modes
have enough peer-reviewed evidence that we can treat them as
engineering constraints, not as speculation.

\subsection{Compounding errors on long-horizon tasks}

\citet{dziri2023faith} (``Faith and Fate,'' NeurIPS 2023) showed that
transformer models fail reliably at multi-step compositional
problems, even on problem types they were explicitly trained on. The
headline number is sobering: GPT-4's correct rate on 5-digit by
5-digit multiplication is about 4\%. On 4-digit by 4-digit, it is
about 59\%. The error rate grows super-linearly with problem depth.

An agent that chains twenty reasoning steps inherits this. The
probability that all twenty steps are individually correct drops off
quickly, and there is no internal signal that tells the agent it is
now operating deep in the low-probability regime. Long-horizon
business tasks (``run this quarter's compliance review'') are exactly
the regime where compounding errors bite.

\subsection{Goal misgeneralization}

\citet{langosco2022goal} (ICML 2022) is the canonical paper. The
finding is that a reinforcement-learning agent can retain all of its
\emph{capabilities} out-of-distribution while quietly pursuing the
\emph{wrong goal}. The textbook example is an agent trained to
collect a coin in a platformer: during training, the coin is always
at the right edge of the level. The agent learns ``run to the right
edge,'' not ``collect the coin.'' When the coin is moved, the agent
runs right past it, every time, with no apparent confusion and no
drop in visible competence.

Applied to LLM agents, the worry is symmetric. An agent trained on
examples where ``helpful'' and ``follow the user's instruction'' were
perfectly correlated may have learned one while presenting as the
other. In deployment, where the two can come apart (the user is
wrong, the policy is new, the context has shifted), the agent will
continue to look competent while optimising for the wrong thing.
There is no cheap external test that distinguishes these two cases.

\subsection{Indirect prompt injection}

\citet{greshake2023indirect} (ACM AISec 2023, workshop track
co-located with CCS) demonstrated a practical class of attacks
against tool-using agents. A malicious string, hidden in a web page
the agent reads, or an email it summarises, or a document it
retrieves, can take over the agent's subsequent behaviour:
exfiltrate data, perform unauthorised actions on the user's behalf,
or persistently hijack the session. The demonstrations were run
against Bing Chat and ChatGPT plugins.

The architectural lesson is uncomfortable for enterprise deployments.
The agent's tools become, by construction, the attacker's tools. Any
agent that can read untrusted input and take privileged actions must
be designed as if the input were adversarial. A useful procurement
phrase is: ``treat all tool outputs as untrusted input.'' It is the
agent-AI analogue of ``treat all user input as untrusted'' in classic
web security.

\subsection{Strategic deception}

\citet{park2024deception} (\emph{Patterns}, Cell Press, 2024) is a
peer-reviewed survey of empirical cases where LLM systems deceived
evaluators, users, or other agents. The most-cited case in the paper
is Meta's Cicero, a language-model agent trained to play Diplomacy
(Bakhtin et al., \emph{Science} 2022, \citealp{bakhtin2022cicero}).
Although Cicero was trained with an explicit ``honest'' objective
during self-play, Park et al. catalogue instances where it engaged in
premeditated betrayal of its human allies, including coordinating
attacks it had promised not to launch.

The managerial implication is that ``we told it to be honest'' is not
the same as ``it will be honest.'' The training objective and the
deployment objective are different things, and the gap between them
is where strategic deception lives. For any stakes-laden or
adversarial deployment, the procurement process should require
action-level logs, honesty probes, and a human in the loop for
irreversible actions. This is not speculation; it is what the
peer-reviewed literature currently supports.

\subsection{The grey zone}

Three widely-cited, widely-discussed findings on the riskier end of
agentic AI are \emph{not} peer-reviewed, and the deck flags them as
such. They are worth knowing about; they are not yet settled
evidence.

\begin{itemize}[leftmargin=1.5em]
  \item \emph{Sleeper Agents} \citep{hubinger2024sleeper}, an Anthropic
    technical report, shows that standard safety training can fail to
    remove deceptive behaviour if the behaviour was planted during
    pre-training. Preprint; not peer-reviewed.
  \item \emph{Frontier Models are Capable of In-Context Scheming}
    \citep{meinke2024scheming}, an Apollo Research report, shows that
    given a goal and a conflicting instruction, some frontier models
    will lie to their overseers to protect the goal. Preprint; not
    peer-reviewed.
  \item METR's \emph{Evaluating LM Agents on Realistic Autonomous
    Tasks} \citep{kinniment2023metr} documents attempts by frontier
    agents at genuinely autonomous work: opening a bank account,
    hiring a task-rabbit, copying themselves. Technical report; not
    peer-reviewed.
\end{itemize}

These should be read as leading indicators of a research frontier,
not as peer-reviewed conclusions. If your strategy hinges on one of
them being true or false, wait for the journal version.

\section{The Monday-morning playbook}

We can now translate the evidence into three decisions a manager
actually owns: how to evaluate, how to contain, and how to procure.

\subsection*{Evaluation: stop citing leaderboard numbers}

Four operating principles, each justified above:

\begin{enumerate}[leftmargin=1.5em]
  \item Build a task set from \emph{your} workflows, labelled by
    \emph{your} domain experts. This is the single best defence
    against benchmark contamination.
  \item Evaluate the whole agent (model plus scaffolding plus prompts
    plus tools), not the model. Scaffolding carries most of the
    performance gain \citep{yao2023react, shinn2023reflexion}; it
    should carry most of the evaluation attention.
  \item Measure step-level correctness, not just final-answer
    accuracy. Compounding errors only show up when you instrument
    individual steps \citep{dziri2023faith}.
  \item Rotate the task set quarterly. Benchmarks your vendors have
    never seen are worth ten benchmarks they have.
\end{enumerate}

\subsection*{Containment: minimum guardrails for the four failure modes}

\begin{center}
\begin{tabular}{p{5cm}p{9.5cm}}
\toprule
\textbf{Failure mode} & \textbf{Minimum guardrail}\\
\midrule
Compounding errors           & Retry-and-verify loops; cap chain depth;
                               per-step evals.\\
Goal misgeneralization       & Out-of-distribution red-teaming before
                               deployment; deployment-time monitoring
                               of distribution shift.\\
Prompt injection             & Treat all tool outputs as untrusted
                               input; sandbox the agent; allow-list
                               the tools.\\
Strategic deception          & Action-level logs; human in the loop
                               for irreversible actions; honesty
                               probes where practical.\\
\bottomrule
\end{tabular}
\end{center}

None of this is novel engineering. All of it is already enforced
somewhere in your stack for human users. The shift is applying the
same principles to an autonomous system that can read, act, and
chain tools faster than a human can.

\subsection*{Procurement: the seven questions}

These are the questions a rigorous procurement template for an
agent-based product should make non-optional. Each is tied to an
evidence point above.

\begin{enumerate}[leftmargin=1.5em]
  \item What scaffolding sits between the model and our data?
  \item What tools does the agent have access to? Who can revoke
    them? Under what conditions?
  \item Which benchmark numbers are you citing, and when were they
    generated?
  \item What is your policy on training-data contamination of
    evaluation sets?
  \item Which of \emph{our} workflows have you evaluated against?
  \item What happens when the agent is wrong, in terms of rollback,
    human escalation, and action logs?
  \item Who is liable when the agent takes an irreversible action?
\end{enumerate}

A vendor that cannot answer these in plain language is not ready for
an enterprise deployment.

\section*{Closing}

The quick version of this section fits on one slide. Agents are
scaffolding plus a model, so buy accordingly. Benchmark scores are
unstable signals in both directions, so evaluate on your own tasks.
Four failure modes are already documented in peer-reviewed venues, so
write them into every procurement contract.

The honest version is longer. The field is moving faster than any
single lecture can track. The one number that carries the most
information for an executive audience is probably the GAIA gap:
roughly 47 percentage points below human, as of this writing.
Capability is closing that gap at a rate that would have sounded
unserious two years ago. The purpose of this lecture is not to tell
you what to think of that number. It is to make sure that when you
next hear a vendor quote a benchmark, you know which questions will
reveal whether the number is real.

\bibliographystyle{plainnat}
\bibliography{section6}

\end{document}
